\documentclass[conference]{IEEEtran}
\IEEEoverridecommandlockouts
\usepackage{cite}
\usepackage{amsmath,amssymb,amsfonts}
\usepackage{array}
\usepackage{graphicx}
\usepackage{textcomp}
\usepackage{xcolor}
\newcommand{\layer}[1]{\textsc{#1}}
\begin{document}

\title{How Is Immune Homeostasis Maintained and Regulated?\\
A Personal Longitudinal Framework for Human Immune Resilience}

\author{\IEEEauthorblockN{Evidence-Grounded Research Proposal}
\IEEEauthorblockA{\textit{Design-only human immune-homeostasis study}}}

\maketitle

\begin{abstract}
Immune homeostasis is commonly described as a balance between protective immunity and damaging inflammation. That description is correct but incomplete: in humans, the relevant balance is individualized, changes across tissues and time, and is influenced by regulatory programs, barrier and microbiota interactions, metabolic context, and exposure-associated immune memory. A population-average cell count or a single blood biomarker therefore cannot define immune homeostasis. This research proposal advances a personal longitudinal immune-homeostasis atlas coupled to bounded human microphysiological feedback testing. The framework estimates a research-state card containing personal baseline distance, tolerance--effector coordination, barrier/microbiota context, memory-context uncertainty, recovery trajectory, data coverage, and abstention. It compares this representation with population-reference and single-score baselines using time-forward evaluation, held-out participants, batch and coverage checks, interval calibration, recovery-time prediction, and layer-attribution audits. A strictly nonclinical human immune--barrier microphysiological platform is proposed only to test preregistered directional feedback hypotheses; it cannot diagnose a person, determine treatment, or replace whole-human validation. The design is motivated by human immune variation, longitudinal multi-omics, regulatory T-cell and mucosal-tolerance evidence, microbiota--immunity crosstalk, and evidence that immune memory leaves durable cell-state traces. No participant has been recruited, no cells have been manipulated, and no results are reported. The intended outcome is an auditable research protocol for determining whether an apparent immune deviation reflects expected adaptation, a persistent dysregulation candidate, or insufficient information.
\end{abstract}

\begin{IEEEkeywords}
immune homeostasis, human immunology, longitudinal multi-omics, regulatory T cells, microbiota, resilience, organ-on-chip, precision medicine
\end{IEEEkeywords}

\section{Introduction}

The immune system must perform two tasks that can appear to conflict. It must recognize and contain threats, but it must also avoid sustained reactivity against self, food, commensal organisms, and healing tissue. Immune homeostasis is the ongoing coordination that makes these tasks compatible. It is not equivalent to weak immunity, low cytokine concentration, or a narrow list of normal cell counts. A robust immune system is able to depart from its baseline when context demands it and to return toward an appropriate operating range when the demand resolves. Failure can occur in either direction: inadequate defense, excessive activation, loss of tolerance, failed resolution, or a response that is appropriately intense but poorly timed for the tissue context.

This definition makes the measurement problem difficult. Human immune states vary greatly among people, even in the absence of overt disease. Brodin and Davis frame this variation as a central feature of human immunology rather than measurement noise \cite{brodin2017}. Tsang \emph{et al.} show substantial baseline variation in healthy people and identify relationships between pre-vaccination immune state and later response \cite{tsang2014}. Tebani \emph{et al.} similarly find high variation across individuals in longitudinal wellness profiling, while selected molecular profiles can be relatively stable within a person \cite{tebani2020}. These observations undermine the premise that an isolated sample can be classified as homeostatic simply because it lies near a population mean.

The problem is also compartmental. Blood provides a valuable but incomplete view of immune state. Mucosal barriers, tissue-resident cells, stromal and epithelial signals, microbial communities, metabolism, circulation, and prior infections all shape the state observed in blood. The microbiota both trains and is constrained by host immunity \cite{belkaid2014,zheng2020,pickard2017}. Regulatory T cells and associated tolerance pathways restrain inappropriate activation, particularly at interfaces rich in dietary and commensal antigen \cite{izcue2006}. Exposure history can leave persistent cellular and epigenetic signatures; long-lasting natural-killer-cell memory provides one example \cite{ruckert2022}. These layers cannot be reduced to a single common axis without losing mechanism and uncertainty.

The practical aspiration is not to replace clinical judgment with an immune score. It is to develop an evidence-ready way of asking a stronger research question: for a particular person, measured context, and sampling coverage, is a coordinated immune trajectory compatible with expected variation and recovery, is it a candidate for further expert investigation, or is it unobservable with the available evidence? This proposal uses the repository's four-agent pipeline to turn that question into a source-bound study design. The Survey Agent limits the claims. The Idea Agent evaluates competing directions. The ExperimentDesign Agent defines variables, comparators, validation, and safeguards. The Author Agent reports the resulting proposal without creating observed outcomes.

\section{Survey Evidence and Research Gap}

\subsection{Homeostasis as regulated resilience}

Population-level studies are indispensable for discovering immune variation, but they do not furnish a universal personal setpoint. In the human variation literature, age, sex, genetic background, exposures, infections, microbiota, and stochastic cell history all contribute to diversity \cite{brodin2017}. In vaccination studies, pre-existing serological state and baseline immune features can influence subsequent responses \cite{tsang2014}. The implication is not that every individual needs an expensive personal assay at every time point. Rather, the intended reference must match the claim. A study claiming to characterize the stability, reactivity, or recovery of one immune system needs repeated measurements and an explicit context model.

Longitudinal multi-omics makes this premise testable. Tebani \emph{et al.} integrate proteomic, transcriptomic, lipidomic, metabolomic, autoantibody, immune-cell, microbiota, and clinical-chemistry information in a wellness cohort \cite{tebani2020}. Their work supports a distinction between high interindividual diversity and more stable intraindividual profiles across selected modalities. The work does not prove that any particular molecular pattern is clinically healthy, nor that blood measurements represent every tissue. It instead supports an individual-baseline design and warns against treating population average as the only standard.

Tolerance and controlled reactivity form a second evidence line. At mucosal interfaces, the immune system must respond to pathogens while tolerating food and commensal antigens. Izcue, Coombes, and Powrie review how regulatory T cells suppress systemic and mucosal activation and discuss TGF-beta, IL-10, CTLA-4, and cell positioning in this regulation \cite{izcue2006}. The review draws from both human and animal work. Its strongest contribution to this proposal is conceptual and mechanistic: tolerance is an active, positioned process. It is not evidence that manipulating a regulatory pathway will safely rebalance a human patient.

The barrier--microbiota layer is equally important. Belkaid and Hand describe microbiota contributions to immune development, containment, regulatory responses, and protective immunity \cite{belkaid2014}. Zheng, Liwinski, and Elinav review multi-directional host--microbiota crosstalk in homeostasis and disease \cite{zheng2020}. Pickard \emph{et al.} connect microbial ecology to pathogen colonization, immune response, and inflammatory-disease context \cite{pickard2017}. These sources argue strongly for including microbial and barrier context in a homeostasis model, but they also warn against a common shortcut: a microbial association does not itself identify a causal target or a generalizable intervention.

Immune history complicates static references further. Ruckert \emph{et al.} map durable clonal and epigenetic features of human adaptive natural killer cells after cytomegalovirus-associated adaptation \cite{ruckert2022}. A durable memory footprint may be a normal consequence of exposure rather than evidence of instability. An appropriate model must therefore represent unobserved or partially observed history as an uncertainty source, not force it into a pathology label.

\begin{table*}[t]
\caption{Survey evidence-to-claim map.}
\label{tab:evidence}
\centering
\renewcommand{\arraystretch}{1.16}
\begin{tabular}{|p{0.22\textwidth}|p{0.36\textwidth}|p{0.31\textwidth}|}
\hline
\textbf{Evidence line} & \textbf{Permitted inference} & \textbf{Prohibited extrapolation} \\
\hline
Human variation and vaccination studies \cite{brodin2017,tsang2014} & Personal baseline and context are necessary for interpreting immune trajectories. & A cohort mean or one measurement diagnoses an individual's immune state. \\
\hline
Longitudinal multi-omics \cite{tebani2020} & Between-person diversity and selected within-person stability can coexist. & Any multi-omic pattern is a validated health or treatment endpoint. \\
\hline
Tolerance and regulation \cite{izcue2006} & Tolerance requires active regulatory pathways and tissue positioning. & Preclinical pathway manipulation is a safe human-rebalancing method. \\
\hline
Microbiota and barrier crosstalk \cite{belkaid2014,zheng2020,pickard2017} & Host--microbe interactions belong in a multi-layer homeostasis model. & A microbial correlation proves causation or identifies a universal therapy. \\
\hline
Human immune memory \cite{ruckert2022} & Exposure history may create durable, cell-state-specific features. & Persistent memory necessarily represents maladaptive dysregulation. \\
\hline
Microphysiological platforms \cite{maharjan2020,kim2016,leung2022} & Human-relevant systems can test selected tissue--immune feedback hypotheses. & A chip recreates whole-body immunity or establishes clinical efficacy. \\
\hline
\end{tabular}
\end{table*}

\subsection{Research gap ledger}

Five linked gaps result from the evidence. \layer{Gap-01} is the static-reference problem: cross-sectional ranges are routinely used where a personal trajectory would be more meaningful. \layer{Gap-02} is the cross-compartment integration problem: blood, barrier, microbiota, tissue, metabolic, and memory information are often collected in isolation. \layer{Gap-03} is the adaptation-versus-dysregulation problem: a transient response, a long-lived memory state, and a persistent pathological process can appear similarly unusual at one time point. \layer{Gap-04} is the human-causal-platform problem: animal and simplified models do not by themselves validate individualized human feedback. \layer{Gap-05} is the rebalancing-endpoint problem: broad immune enhancement or suppression lacks a validated person-specific endpoint for restored contextual regulation.

The gaps are mutually reinforcing. A one-number score can conceal that a high inflammatory module reflects a short-lived adaptation, a missing tissue compartment, a batch shift, or a true persistent deviation. A more complex model is not automatically safer: without a coverage record, a defined recovery horizon, and an abstention option, it can merely hide greater uncertainty behind more features.

\section{Idea Agent: Direction Search and Selection}

The Idea Agent examined four directions. A cross-sectional immune-health atlas is scientifically valuable and remains a comparator, but it cannot infer a person's recovery or distinguish stable individual variation from deviation. A one-biomarker blood score was rejected because it collapses tissue, memory, tolerance, and exposure context. A broad modulation direction was rejected because it mistakes score movement for an individual-specific rebalancing endpoint and would create unacceptable medical overclaiming.

The selected direction is a personal longitudinal immune-homeostasis atlas coupled to bounded human microphysiological counterfactual testing. The model does not assert that every homeostatic layer is observable. It represents what is measured, what is unknown, and whether the data can support a research-state interpretation. Its novelty is the connection of four elements that are normally separated: personal operating range, layer-specific attribution, a recovery trajectory, and a narrowly scoped platform test for a preregistered feedback relation.

\begin{figure*}[t]
\centering
\fbox{\begin{minipage}{0.93\textwidth}
\centering\small
\textbf{Survey Agent} \quad $\longrightarrow$ \quad \textbf{Idea Agent} \quad $\longrightarrow$ \quad \textbf{ExperimentDesign Agent} \quad $\longrightarrow$ \quad \textbf{Author Agent} \\
\vspace{0.35em}
Human evidence map and gap ledger \quad $\longrightarrow$ \quad competing directions and selected D2 \quad $\longrightarrow$ \quad trajectories, comparators, safeguards, and outcomes \quad $\longrightarrow$ \quad source-bound IEEE proposal \\
\vspace{0.85em}
\fbox{\parbox{0.83\textwidth}{\centering\textbf{Selected research target:} a personal, context-qualified operating range across \emph{tolerance--effector coordination}, \emph{barrier/microbiota context}, and \emph{exposure-memory context}; a deviation is evaluated as recovery, persistent candidate, or \textbf{abstain}, never as an automated diagnosis.}}
\end{minipage}}
\caption{Traceable four-agent workflow and selected immune-homeostasis target. The pipeline creates a proposal and reports no new observed data.}
\label{fig:workflow}
\end{figure*}

The central hypothesis is deliberately falsifiable. If a personal trajectory does not improve interval calibration, recovery prediction, or layer attribution over a population reference, the additional complexity is not justified. If tissue and data limitations dominate, the model must abstain. If the human microphysiological platform does not support the preregistered direction of a proposed feedback relation, that relation should not be promoted as a mechanistic explanation. The framework earns complexity only if it increases transparent and bounded knowledge.

\section{Research Questions and Formal Hypotheses}

\textbf{RQ1:} Can repeated, individual-level measurements estimate an operating range more honestly than a population reference? \textbf{H1} predicts that a personal model will better calibrate expected within-person fluctuation and recovery time on later measurement windows.

\textbf{RQ2:} Can a multi-layer representation make deviations more interpretable? \textbf{H2} predicts that an audit will more consistently attribute a state to tolerance--effector, barrier/microbial, memory-context, or coverage uncertainty than a pooled anomaly score.

\textbf{RQ3:} Can a human-relevant nonclinical platform test a limited feedback claim? \textbf{H3} predicts that a preregistered immune--barrier relation can be assessed for directional consistency in a bounded microphysiological system. The hypothesis does not predict clinical outcome, whole-body response, or therapeutic effect.

\textbf{RQ4:} When is an immune-homeostasis state not estimable? \textbf{H4} predicts that missing tissue access, insufficient temporal coverage, uncertain history, or consent limitations will identify cases where abstention is more responsible than a forced score.

For person $i$ at time $t$, let the research-state vector be
\begin{equation}
\mathbf{h}_{i,t}=\left[b_{i,t},q_{i,t},m_{i,t},r_{i,t},c_{i,t}\right],
\label{eq:state}
\end{equation}
where $b$ is personal-baseline distance, $q$ denotes tolerance--effector coordination, $m$ is barrier/microbiota and memory context, $r$ is recovery-trajectory evidence, and $c$ is coverage and consent quality. The vector is not a biological truth claim. It is a structured accounting of the available evidence and uncertainty.

The proposal estimates an individual-specific interval rather than a point norm:
\begin{equation}
\mathcal{I}_{i,t+h}=f\left(\mathbf{h}_{i,\leq t},X_{i,\leq t},C_i\right),
\label{eq:interval}
\end{equation}
where $X$ contains only consent-authorized observations available by $t$, $C_i$ encodes coverage and consent constraints, and $h$ is a preregistered follow-up horizon. The output must include the option $\varnothing$ for non-estimability. A model that produces a narrow interval when relevant tissue, temporal, or consent information is absent is miscalibrated by design.

\section{ExperimentDesign Agent: Study Design}

\subsection{Execution boundary and data architecture}

This is a \textbf{design-only} protocol. It does not recruit participants, collect samples, assign a vaccine or treatment, engineer cells, expose cells to pathogens, or make a medical decision. If later implemented, the observational component would use consented, de-identified research data under ethics, privacy, biosafety, and clinical-governance approvals. The platform component would use only approved nonclinical materials and institutional procedures. No output is disclosed as a diagnosis, a health recommendation, or an indication to modify therapy.

The unit of analysis is a consented participant-by-measurement window. Data are represented through versioned source cards. A card describes the modality, collection time, tissue or compartment scope, permitted use, processing batch, missingness, delay, consent status, and known confounders. Potential modality families include aggregated immune-cell phenotypes, targeted protein summaries, transcriptomic or chromatin summaries, metabolite indicators, noninvasive microbial features, and authorized contextual metadata. The specific assays and sampling schedules remain choices for an approved future protocol; this report deliberately does not supply laboratory instructions.

The state card contains six fields. Personal baseline distance reports whether a measurement lies inside or outside an estimated individual interval. Tolerance--effector coordination reports predefined regulatory and activation modules separately. Barrier/microbiota context records authorized mucosal and microbial proxies, not a causal microbial diagnosis. Memory context records the degree to which known or unknown exposure history may explain a durable state. Recovery evidence describes the trajectory relative to a preregistered horizon. Coverage reports tissue, time, quality, consent, and batch limitations. An abstention output overrides the other fields whenever a defensible interpretation is not possible.

\begin{table*}[t]
\caption{Homeostasis state variables, outputs, and safeguards.}
\label{tab:variables}
\centering
\renewcommand{\arraystretch}{1.13}
\begin{tabular}{|p{0.18\textwidth}|p{0.29\textwidth}|p{0.25\textwidth}|p{0.17\textwidth}|}
\hline
\textbf{Construct} & \textbf{Authorized representation} & \textbf{Research output} & \textbf{Safeguard} \\
\hline
Personal baseline & Repeated de-identified cellular and molecular summaries & Operating-range interval and baseline distance & Not a diagnostic reference range. \\
\hline
Tolerance--effector coordination & Predefined aggregate regulatory and activation modules & Coordination trajectory plus uncertainty & A module does not stand for a complete cell function. \\
\hline
Barrier/microbiota context & Authorized noninvasive microbial, metabolite, and barrier-associated summaries & Context flag and layer attribution & Association is not a causal target. \\
\hline
Memory context & Consent-authorized exposure information and memory-state summaries & History uncertainty adjustment & No inference about unreported exposure. \\
\hline
Platform feedback readout & Predefined aggregate immune--barrier measures in an approved nonclinical system & Directional-consistency test & No clinical, whole-body, or efficacy claim. \\
\hline
Coverage and consent & Tissue scope, data age, missingness, batch, permissions & Reliability card and abstention & Missing data cannot become normality. \\
\hline
\end{tabular}
\end{table*}

\subsection{Comparators and bounded platform role}

The first comparator is a cross-sectional population-reference model using only prespecified and ethically justified stratification variables. It establishes the value of existing reference practice. The second is a pooled multi-omic anomaly score without layer attribution. It tests whether additional structure yields actual interpretive value rather than a more elaborate number. The third is the personal state-space proposal, which produces individual intervals, layer outputs, recovery evidence, coverage flags, and abstention. The fourth is not a competing predictor: it is a bounded platform comparator for one predeclared immune--barrier feedback hypothesis.

Human microphysiological systems provide a promising but limited platform. Immunocompetent organ-on-chip models can represent selected tissue architecture, immune components, and microenvironmental interactions \cite{maharjan2020}. Human gut-on-chip work illustrates that barrier, microbial, mechanical, and inflammatory contributions can be separated in a controlled system \cite{kim2016}. Design guidance for organ-on-chip systems emphasizes fit-for-purpose choices and remaining translation limits \cite{leung2022}. Accordingly, the proposed platform is not called a human immune system in miniature. It is a narrow testbed: before platform work begins, investigators specify one feedback hypothesis, the aggregate readouts that would be directionally consistent or inconsistent with it, the permissible scope, and the claim that will \emph{not} be made from the result.

\section{Evaluation, Interpretation, and Failure Analysis}

\subsection{Deployment-realistic evaluation}

The study starts with a protocol freeze. Consent language, data dictionary, window definitions, permitted metadata, homeostatic modules, personal-baseline method, recovery horizon, missingness rules, primary metrics, and communication controls are fixed before model fitting. Later changes are versioned as a new protocol. This prevents a model from defining homeostasis only after it sees which definition performs best.

Evaluation uses time-forward holdouts within people, so a future measurement is never used to estimate an earlier state. Held-out participants test whether the framework's uncertainty and abstention behavior transport, even though the expected personal baseline itself cannot be copied from another person. Where provenance permits, sites or processing batches are held out. The study also simulates incomplete coverage and realistic measurement latency. This is essential because an apparently confident immune score may be driven by data-rich cohorts, stable batches, or repeated access to a compartment not available in later use.

Primary measures are personal-interval calibration, recovery-time calibration, coverage of uncertainty intervals, attribution consistency, abstention appropriateness, and coverage-stratified disparity. Classification accuracy is deliberately secondary: a model that detects deviation but cannot state whether it reflects ordinary variation, recovery, missing tissue, or persistent candidate dysregulation cannot responsibly support a research interpretation. A structured audit samples cards from all outcome types and asks reviewers whether the time cut-off, source cards, module attribution, interval width, and abstention logic are internally consistent.

\begin{figure}[t]
\centering
\fbox{\begin{minipage}{0.94\columnwidth}
\centering\footnotesize
\textbf{1. Personal baseline}\\
Repeated, consent-limited observations\\
$\Downarrow$ (interval and coverage)\\
\textbf{2. Regulatory coordination}\\
Tolerance and effector modules kept distinct\\
$\Downarrow$ (barrier and memory context)\\
\textbf{3. Resilience trajectory}\\
Expected fluctuation, recovery, or persistent candidate\\
$\Downarrow$ (tissue, consent, and data quality)\\
\textbf{4. Research-state interpretation}\\
\fbox{stable} \quad \fbox{resolved adaptation} \quad \fbox{candidate} \quad \fbox{abstain}
\end{minipage}}
\caption{Personal homeostasis decision path. Each stage retains uncertainty and scope limitations; no state is a diagnosis or treatment recommendation.}
\label{fig:decisionpath}
\end{figure}

\subsection{Outcome branches and platform audit}

The protocol permits four conclusion labels. \emph{Individual baseline stable} requires calibration of the personal operating range and a fluctuation pattern consistent with that range on later windows. \emph{Contextual adaptation resolved} requires an adequately covered deviation that returns according to the preregistered recovery structure. \emph{Persistent dysregulation candidate} requires a persistent, layer-attributed pattern, adequate coverage, and human reviewer agreement; it authorizes only further research review, not diagnosis. \emph{Model or data abstain} is required when tissue coverage, temporal evidence, consent, calibration, or mechanism resolution is inadequate.

The platform sub-study has an equally narrow audit. A computational model may nominate a feedback hypothesis only from a predeclared set. Before platform data are examined, the protocol identifies the model's directional prediction, the aggregate readouts, and the conditions under which the relation would be judged unsupported. A result that is directionally consistent can support further mechanistic study; it cannot validate the participant model, establish a clinical pathway, or prove that a future intervention will rebalance immunity. An inconsistent result is valuable because it forces revision of the model's causal interpretation rather than a hidden post hoc story.

\begin{table*}[t]
\caption{Predeclared evaluation scenarios and maximum permissible interpretations.}
\label{tab:scenarios}
\centering
\renewcommand{\arraystretch}{1.13}
\begin{tabular}{|p{0.21\textwidth}|p{0.27\textwidth}|p{0.22\textwidth}|p{0.17\textwidth}|}
\hline
\textbf{Scenario} & \textbf{Expected state-card behavior} & \textbf{Evaluation question} & \textbf{Maximum permissible interpretation} \\
\hline
Stable personal pattern, unusual population value & Personal interval remains calibrated; population difference is visible but not dominant & Does the model avoid pathologizing stable individual variation? & Consistent with a personal research baseline. \\
\hline
Transient regulatory/effector deviation with recovery & Layer attribution and recovery interval are shown; memory context retained & Does the model distinguish a resolved adaptation from a persistent candidate? & Consistent with resolved contextual adaptation. \\
\hline
Persistent deviation with adequate multi-layer coverage & Reproducible layer pattern, uncertainty, and reviewer audit are available & Does evidence justify escalation to research expert review? & Persistent dysregulation candidate; never diagnosis. \\
\hline
Sparse tissue or temporal coverage & Wide interval, tissue/consent limitations, and abstention override & Does absence of measurement lower confidence rather than create false normality? & Abstain from a homeostasis-state interpretation. \\
\hline
Platform inconsistency & Directional prediction and predefined readout disagree & Does the proposal revise the causal explanation rather than reinterpret the result? & Feedback hypothesis unsupported in this model scope. \\
\hline
\end{tabular}
\end{table*}

Failure analysis is a primary deliverable. The personal model may not beat a population reference after the cost of additional sampling and computation is considered. The state modules may be unstable across batches. Barrier and tissue uncertainty may be too large to support an interpretable interval. A memory state may confound the apparent recovery path. The platform might reject a feedback hypothesis or reveal that its result depends on the system's limited architecture. Each failure produces a specific response: simplify the model, improve provenance, narrow the claim, collect better longitudinal evidence under a new approved protocol, or abstain. None should be repaired by silently defining the participant as dysregulated.

\section{Prospective Research Audit and Interpretation Rules}

Retrospective datasets are necessary for model development but cannot demonstrate whether the framework would remain honest when measurements arrive sequentially, when a sample is missing, or when a contextual event is only partly recorded. The proposal therefore specifies a future prospective \emph{research shadow phase}. During this phase, authorized de-identified measurements would be processed according to the frozen protocol and converted into a state card. The card would be available only to an approved research review group. It would not be returned to participants, added to a health record, displayed as medical advice, or used to change care. Later measurements would be used to assess calibration and recovery predictions, but they could not be inserted retrospectively into the earlier card.

The shadow phase distinguishes a homeostasis model from a descriptive dashboard. Each issuance records the exact information cut-off, source-card versions, consent scope, baseline interval, layer-specific outputs, uncertainty, coverage status, preliminary research label, and reviewer rationale. This record enables a later evaluator to ask what the model actually knew at the relevant time. If a late measurement, an amended exposure history, or a processing correction would have altered the conclusion, the difference becomes part of the evaluation rather than disappearing in a cleaned final dataset.

The review group has distinct roles. A computational reviewer checks temporal leakage, provenance, batch handling, interval construction, and abstention logic. An immunology reviewer checks whether layer attribution is biologically coherent within the proposal's limits. A governance reviewer checks consent, data-minimization, communication, and potential stigmatization risks. None of these roles turns a research label into a diagnosis. Disagreement is preserved as an auditable finding: persistent reviewer disagreement may itself indicate that the state vector lacks sufficient biological resolution or that the claim is too broad.

\begin{table}[t]
\caption{Predeclared research shadow-audit record.}
\label{tab:audit}
\centering
\renewcommand{\arraystretch}{1.12}
\begin{tabular}{|p{0.27\columnwidth}|p{0.61\columnwidth}|}
\hline
\textbf{Record field} & \textbf{Required content} \\
\hline
Information cut-off & Measurement time, source versions, authorized metadata, and later information excluded from the state card. \\
\hline
Coverage and consent & Tissue scope, modality completeness, batch status, access restriction, and reason for any abstention. \\
\hline
State representation & Personal interval, layer outputs, memory-context uncertainty, recovery horizon, and interval width. \\
\hline
Interpretation path & Research label, named limitations, maximum permissible claim, and reasons alternate labels were not selected. \\
\hline
Reviewer audit & Computational, immunology, and governance checks; disagreements and protocol deviations retained. \\
\hline
Later evaluation & Predeclared later window, calibration contribution, recovery evidence, and model-version identity. \\
\hline
\end{tabular}
\end{table}

Interpretation is governed by asymmetric error costs. Calling ordinary individual variation abnormal can create anxiety, stigma, unnecessary investigation, and inappropriate pressure to alter immunity. Treating a persistent dysregulation candidate as ordinary variation can obscure a research signal. The model cannot resolve this trade-off by threshold alone. Before a shadow phase, the research team must specify which errors matter for its stated scientific purpose, which claims are permitted, and when uncertainty requires abstention. The proposal does not endorse a common threshold across laboratories or populations; the minimum common requirement is transparent reporting of the threshold rationale and the evidence it needs.

The same discipline applies to the microphysiological platform. A platform result must be reported with its model scope: which human cell or tissue features are represented, which are not represented, which readouts are aggregate, and which environmental, endocrine, vascular, neural, or systemic feedbacks remain outside the platform. A directional result may strengthen the plausibility of a narrowly framed relation between an immune and barrier feature. It cannot establish that the relation drives a person's in-vivo trajectory, that it is sufficient to cause disease, or that manipulating it would restore health. This limit is a methodological asset. It prevents an observational association from acquiring causal authority merely because it was reproduced in a simplified setting.

Model updates require equally strict discipline. An update may be triggered by a prespecified calibration failure, a source retirement, a material batch shift, a consent-policy change, or a revised measurement definition. It may not be triggered solely because a recently inspected subset made a revised model look more convincing. Each version retains its own source cards, state definitions, validation period, audit records, and output rules. The result is a lineage of claims: reviewers can determine whether an apparent improvement reflects better calibration, different data, narrower coverage, or an altered interpretation policy.

Finally, the protocol asks a red-team style question at each milestone: could a stable individual signature be mistaken for pathology, could a missing tissue be mistaken for absence of dysregulation, could prior immune memory be mistaken for a current process, could batch correction hide a true temporal change, or could a chip result be communicated as a treatment signal? The review does not attempt clinical intervention and does not seek operational methods for changing immune function. Its function is to test whether the proposal's own limits survive realistic analytical and interpretive pressure.

\section{Governance, Equity, and Limitations}

Human immune information is sensitive because it can be misread as a diagnosis, a prediction of disease, or a justification for treatment. The proposal requires purpose-limited consent, de-identification, data minimization, versioned access control, and clear separation between research labels and clinical care. A participant or clinician must not receive an algorithmic homeostasis score as advice. Any future return of research findings requires its own ethics, analytical-validity, clinical-validity, and communication plan.

Equity matters at two levels. First, a model trained in one population may treat unrepresented biological variation as deviation. Second, individuals with less frequent sampling, unavailable tissues, or different exposure documentation can receive falsely narrow estimates. The state card must show coverage and consent limitations as prominently as the estimated layers. Abstention is not a solution to unequal evidence, but it prevents the misleading claim that sparse observation demonstrates immune equilibrium. A responsible program would use abstention patterns to guide data-governance and research-investment decisions, not to assign biological deficiency to people or communities.

The limitations are substantial. Blood is not tissue. Microbiota features are context dependent. Omics measurements contain technical and batch variation. Homeostatic modules are modelling choices, not natural laws. A durable immune-memory signature can be normal. Human microphysiological systems provide selected controlled contexts, not complete systemic physiology. The framework does not eliminate these limitations; it makes their influence visible, measures whether predictions survive them, and prohibits the model from speaking when the evidence does not support it.

\section{Conclusion}

Immune homeostasis is maintained by active, multi-layer regulation rather than by a universal static immune value. Human variation, longitudinal stability, tolerance circuits, barrier--microbiota crosstalk, and exposure-associated memory make a one-score account scientifically inadequate. The practical research opportunity is to define homeostasis as context-qualified resilience: a coordinated system can change, recover, and remain interpretable within the limits of its measurement coverage.

This four-agent proposal operationalizes that opportunity. The Survey Agent establishes a human-evidence boundary and a gap ledger. The Idea Agent rejects the universal-score and broad-modulation narratives in favor of a personal trajectory. The ExperimentDesign Agent specifies comparators, time-forward evaluation, abstention, and a bounded human-platform mechanism test. The Author Agent preserves the distinction between proposal and result. The decisive test is not whether the system produces a more impressive immune number; it is whether it can give a transparent reason to regard a measured trajectory as expected variation, resolved adaptation, persistent research candidate, or insufficient evidence. If it cannot, it should abstain. If it can, it offers a credible foundation for the future science of immune rebalancing without claiming that rebalancing has already been achieved.

\begin{thebibliography}{00}
\bibitem{brodin2017} P. Brodin and M. M. Davis, ``Human immune system variation,'' \emph{Nature Reviews Immunology}, vol. 17, pp. 21--29, 2017, doi: 10.1038/nri.2016.125.
\bibitem{tsang2014} J. S. Tsang \emph{et al.}, ``Global analyses of human immune variation reveal baseline predictors of postvaccination responses,'' \emph{Cell}, vol. 157, no. 2, pp. 499--513, 2014, doi: 10.1016/j.cell.2014.03.031.
\bibitem{tebani2020} A. Tebani \emph{et al.}, ``Integration of molecular profiles in a longitudinal wellness profiling cohort,'' \emph{Nature Communications}, vol. 11, art. 4487, 2020, doi: 10.1038/s41467-020-18148-7.
\bibitem{belkaid2014} Y. Belkaid and T. W. Hand, ``Role of the microbiota in immunity and inflammation,'' \emph{Cell}, vol. 157, no. 1, pp. 121--141, 2014, doi: 10.1016/j.cell.2014.03.011.
\bibitem{zheng2020} D. Zheng, T. Liwinski, and E. Elinav, ``Interaction between microbiota and immunity in health and disease,'' \emph{Cell Research}, vol. 30, pp. 492--506, 2020, doi: 10.1038/s41422-020-0332-7.
\bibitem{pickard2017} J. M. Pickard, M. Y. Zeng, R. Caruso, and G. Nunez, ``Gut microbiota: Role in pathogen colonization, immune responses, and inflammatory disease,'' \emph{Immunological Reviews}, vol. 279, no. 1, pp. 70--89, 2017, doi: 10.1111/imr.12567.
\bibitem{izcue2006} A. Izcue, J. L. Coombes, and F. Powrie, ``Regulatory T cells suppress systemic and mucosal immune activation to control intestinal inflammation,'' \emph{Immunological Reviews}, vol. 212, pp. 256--271, 2006, doi: 10.1111/j.0105-2896.2006.00423.x.
\bibitem{ruckert2022} T. Ruckert \emph{et al.}, ``Clonal expansion and epigenetic inheritance of long-lasting NK cell memory,'' \emph{Nature Immunology}, vol. 23, pp. 1551--1563, 2022, doi: 10.1038/s41590-022-01327-7.
\bibitem{maharjan2020} S. Maharjan, B. Cecen, and Y. S. Zhang, ``3D immunocompetent organ-on-a-chip models,'' \emph{Small Methods}, vol. 4, art. 2000235, 2020, doi: 10.1002/smtd.202000235.
\bibitem{kim2016} H. J. Kim, H. Li, J. J. Collins, and D. E. Ingber, ``Contributions of microbiome and mechanical deformation to intestinal bacterial overgrowth and inflammation in a human gut-on-a-chip,'' \emph{Proceedings of the National Academy of Sciences}, vol. 113, no. 1, pp. E7--E15, 2016, doi: 10.1073/pnas.1522193112.
\bibitem{leung2022} C. M. Leung \emph{et al.}, ``A guide to the organ-on-a-chip,'' \emph{Nature Reviews Methods Primers}, vol. 2, 2022, doi: 10.1038/s43586-022-00118-6.
\end{thebibliography}

\end{document}
